\documentclass[10pt,a4paper]{amsart}
\usepackage[margin=19mm]{geometry}
\usepackage{amsmath,amssymb,mathtools}
\usepackage[scr=rsfs,cal=euler]{mathalfa}
\usepackage{xfrac}
\usepackage[unicode,hidelinks,bookmarks=false]{hyperref}
\newcommand{\ie}{i.e.\ }
\newcommand{\cf}{cf.\ }
\newcommand{\resp}{respectively\ }
\newcommand{\Subsection}{\subsection}
\providecommand{\nobreakdash}{\nobreak\hbox{-}}
\setlength{\emergencystretch}{2em}
\multlinegap0pt
\usepackage[shortlabels]{enumitem}
\usepackage{amssymb}
\usepackage{mathtools}
\usepackage{bm}
\newtheorem{definition}{Definition}
\newtheorem{proposition}{Proposition}
\newtheorem{lemma}{Lemma}
\newcommand{\CC}{\mathbb{C}}
\title{Local Rank-One Logarithmic Instability for the Mixed Hessian of the Dispersionless Toda $\tau$-Function}
\author{Oleg Alekseev}
\date{Source: arXiv:2604.00240v2 (CC BY 4.0)}
\begin{document}
\maketitle

\noindent\emph{Source and licence.} Local Rank-One Logarithmic Instability for the Mixed Hessian of the Dispersionless Toda $\tau$-Function. Version arXiv:2604.00240v2 (CC BY 4.0), distributed by arXiv under the Creative Commons Attribution 4.0 International licence, http://creativecommons.org/licenses/by/4.0/, as recorded in the arXiv licence metadata for this exact version and shown on the abstract page https://arxiv.org/abs/2604.00240v2. Full licence notice: \texttt{licenses/arxiv-2604.00240-CC-BY-4.0.txt}.\par\smallskip

\noindent\textit{Editorial scope.} Counted unit: the article's lemma lem:critical-coeff-asymp (source0.tex lines 3098--3130). The article's complete original proof of it is delivered with it (source0.tex lines 3132--3190, one \texttt{proof} environment of 59 lines). No mathematics is written, repaired or re-derived: the statement and the proof are the article's own lines, in the article's own notation, and no retained line is substituted or re-flowed.  Retained uncounted as the source-local context the proof needs: lines 470--472; lines 506--512; lines 728--732; lines 735--775; lines 761--763; lines 765--767; lines 791--802; lines 829--841; lines 902--925; lines 964--969; lines 984--994; lines 1481--1587; lines 3021--3039; lines 3031--3034.  Nothing was omitted from any retained span, so the package carries no omission record.  No retained line contains a citation, so the wrapper carries no bibliography and no bibliography entry.  No retained line contains a figure environment, an image-inclusion command or drawing code, so no image is delivered and the package declares no asset dependency.  Editorial formatting only, in addition: an AMS \texttt{amsart} preamble that restates the article's own declarations and macros used by the retained lines, namely the environments \texttt{definition}, \texttt{proposition}, \texttt{lemma} on separate counters, together with the standard packages those lines need.  The retained environments print in this document as Definition 1, Proposition 1, Lemma 1 in order of appearance, whereas the article prints the same items with the numbering of its own full document.  The complete native source of the article as distributed in the arXiv e-print for this exact version is bundled unchanged as \texttt{sources/197577/source0.tex}.
	\begin{equation}\label{eq:inv-eq}
		U(x;\zeta)=1+\sum_{n=1}^N \zeta_n\,x^{s_n}\,U(x;\zeta)^{s_n}.
	\end{equation}

\begin{definition}[Radius of analyticity]\label{def:rho-star}
	Let $\rho_*(\zeta)$ denote the radius of the largest disk centered at $x=0$ on which the Taylor
	branch of \eqref{eq:inv-eq} is analytic:
	\begin{equation}\label{eq:rho-star}
		\rho_*(\zeta):=\sup\{\rho>0:\ U(\cdot;\zeta)\ \text{analytic on }|x|<\rho\}.
	\end{equation}
\end{definition}

\begin{equation}\label{eq:branch-eq}
	F(U(x;\zeta),x;\zeta)=0,
	\qquad
	U(0;\zeta)=1.
\end{equation}

\begin{proposition}[Branch points and the characteristic system]\label{prop:xi-star-system}
	Fix $\zeta$ and consider the Taylor branch $x\mapsto U(x;\zeta)$ determined by \eqref{eq:branch-eq}.
	Assume that this branch admits analytic continuation in a slit neighborhood of a point $x_*\in\CC$
	with $0<|x_*|<\infty$, and that $x_*$ is a singular point on the Taylor sheet. Let $\lambda$ denote
	the corresponding finite branch value attained by that continuation at $x_*$. Then $(x_*,\lambda)$
	satisfies the \emph{characteristic system}
	\begin{equation}\label{eq:char-system-F}
		F(\lambda,x_*;\zeta)=0,
		\qquad
		\partial_yF(\lambda,x_*;\zeta)=0.
	\end{equation}
	Equivalently,
	\begin{equation}\label{eq:char-system}
		\lambda
		=
		1+\sum_{n=1}^N \zeta_n\,x_*^{s_n} \lambda^{s_n},
		\qquad
		1
		=
		\sum_{n=1}^N s_n\,\zeta_n\,x_*^{s_n} \lambda^{s_n-1}.
	\end{equation}

	Conversely, suppose that the Taylor branch admits analytic continuation to a
	punctured neighborhood of $x_*$ with branch value $\lambda$, that
	$(x_*,\lambda)$ satisfies \eqref{eq:char-system-F}, and that the
	nondegeneracy conditions
	\begin{equation}\label{eq:Fx-nondeg}
		\partial_xF(\lambda,x_*;\zeta)\neq0,
	\end{equation}
	and
	\begin{equation}\label{eq:Fyy-nondeg}
		\partial_y^2F(\lambda,x_*;\zeta)\neq0
	\end{equation}
	hold. Then $x_*$ is a square--root branch point of the Taylor sheet: in a neighborhood of $x_*$ one has a
	local Puiseux expansion
	\[
		U(x;\zeta)=\lambda + c\,\sqrt{x-x_*}+O(x-x_*),
		\qquad c\neq0,
	\]
	for a suitable choice of the square--root branch.
\end{proposition}

	\begin{equation}
		\partial_xF(\lambda,x_*;\zeta)\neq0,
	\end{equation}

	\begin{equation}
		\partial_y^2F(\lambda,x_*;\zeta)\neq0
	\end{equation}

\begin{definition}[Dominant $s$--orbit at the analytic boundary]\label{def:dominant-orbit}
	Fix $\zeta$ and assume the Taylor branch is analytic in $|x|<\rho_*(\zeta)$, where $\rho_*(\zeta)$ is
	as in Definition~\ref{def:rho-star}. Let $s:=\gcd(s_1,\dots,s_N)$ so that
	$U(e^{2\pi i/s}x;\zeta)=U(x;\zeta)$. We say that the analytic boundary is \emph{dominant modulo the $s$--symmetry} if the chosen Taylor sheet has
	exactly $s$ singular points on the circle $|x|=\rho_*(\zeta)$, forming a single orbit
	\[
		\bigl\{e^{2\pi i j/s}\,x_*(\zeta)\bigr\}_{j=0}^{s-1},
		\qquad |x_*(\zeta)|=\rho_*(\zeta),
	\]
	and no other singularities on that circle. We refer to $x_*(\zeta)$ as a chosen representative of the
	dominant orbit.
\end{definition}

\begin{definition}[Simple critical point]\label{def:simple-critical}
	A parameter $\zeta_c\in\mathcal C$ is called a \emph{simple critical point}
	if, for one representative $x_{*,c}$ of the dominant
	orbit and the corresponding finite branch value $\lambda_c$ on the chosen Taylor
	sheet, one has
	\[
		F(\lambda_c,x_{*,c};\zeta_c)=0,
		\qquad
		\partial_yF(\lambda_c,x_{*,c};\zeta_c)=0,
		\qquad
		\partial_y^2F(\lambda_c,x_{*,c};\zeta_c)\neq0.
	\]
\end{definition}

\begin{lemma}[Local continuation of the critical branch parameters]
\label{lem:local-branch-data}
Let \(\zeta_c\in\mathcal C\) be a simple critical point, and choose a
representative critical point \(x_{*,c}\) together with the corresponding
branch value \(\lambda_c\) on the chosen Taylor sheet. Then there exist a
neighborhood \(\mathcal O\) of \(\zeta_c\) and unique \(C^1\) maps
\[
\zeta\mapsto x_*(\zeta),\qquad
\zeta\mapsto \lambda(\zeta),\qquad
\zeta\in\mathcal O,
\]
such that \((x_*(\zeta_c),\lambda(\zeta_c))=(x_{*,c},\lambda_c)\) and
\[
F(\lambda(\zeta),x_*(\zeta);\zeta)=0,
\qquad
\partial_yF(\lambda(\zeta),x_*(\zeta);\zeta)=0.
\]
If, in addition, the continued point \(x_*(\zeta)\) remains a representative
of the unique dominant orbit for \(\zeta\in\mathcal O\), then
\[
\rho_*(\zeta)=|x_*(\zeta)|
\]
on \(\mathcal O\), and hence \(\rho_*\) is \(C^1\) on \(\mathcal O\).
\end{lemma}

\begin{equation}\label{eq:lambda-def}
	\lambda(\varepsilon)
	:=
	\lim_{x\to x_*(\varepsilon)}
	U(x;\zeta(\varepsilon)),
\end{equation}

\begin{definition}[$\Delta$--domains and $\Delta$--analyticity]\label{def:delta-domain}
	Fix a point $\xi_*\in\CC\setminus\{0\}$ and an opening angle $\phi\in(0,\pi/2)$. A \emph{$\Delta$--domain at $\xi_*$} is a domain
	of the form
	\[
		\Delta(\xi_*,\phi)
		:=
		\Bigl\{x\in\CC:\ |x|<|\xi_*|+\eta,\ x\neq\xi_*,\ |\arg(\xi_*-x)|>\phi\Bigr\},
	\]
	for some thickness parameter $\eta>0$. A function $G$ is \emph{$\Delta$--analytic at $\xi_*$} if it admits an analytic
	continuation to $\Delta(\xi_*,\phi)$ for some $\phi\in(0,\pi/2)$.
\end{definition}

\begin{definition}[Uniform continuation hypothesis near the dominant orbit]\label{def:uniform-delta}
	Let $\zeta(\varepsilon)$ be a $C^1$ path with $\zeta(0)=\zeta_c\in\mathcal C$ a simple critical point
	(Definition~\ref{def:simple-critical}). Assume that for $0\le\varepsilon\le\varepsilon_1$ the Taylor branch
	$U(\,\cdot\,;\zeta(\varepsilon))$ lies in the dominant regime of Definition~\ref{def:dominant-orbit} with a unique
	dominant $s$--orbit of simple branch points.
	Write $s=\gcd(s_1,\dots,s_N)$, set $z:=x^{s}$, and, for notational convenience, denote by
	$\widehat U(z;\varepsilon)$ the same Taylor sheet expressed in the variable $z$.
	Let $x_*(\varepsilon)$ be a representative dominant singularity and set $z_*(\varepsilon):=x_*(\varepsilon)^s$.
	Let $\lambda(\varepsilon)$ denote the corresponding branch value of the chosen Taylor sheet at $z_*(\varepsilon)$, that is,
	\[
		\lambda(\varepsilon):=
		\lim_{z\to z_*(\varepsilon)}\widehat U(z;\varepsilon),
	\]
	where $z$ is on the chosen Taylor sheet.
	We say that the path satisfies a \emph{uniform square-root continuation
condition} if there exist constants \(\varepsilon_0>0\),
\(\phi\in(0,\pi/2)\), \(\eta_0>0\), and \(B_0,C_0,c_0>0\) such that, for every
\(0<\varepsilon<\varepsilon_0\), the following properties hold:
	\begin{enumerate}[label=(U\arabic*),leftmargin=2.5em]
		    \item\label{U1} \emph{Uniform indented-domain continuation on an indented disk.}
		      Set
		      \[
			      R(\varepsilon):=|z_*(\varepsilon)|+\eta_0/2,
			      \qquad
			      \Omega_\varepsilon:=\{\,|z|<R(\varepsilon)\,\}\cap \Delta(z_*(\varepsilon),\phi),
		      \]
			    where \(\Delta(z_*(\varepsilon),\phi)\) is the indented domain from
		      Definition~\ref{def:delta-domain}. The Taylor sheet
		      \(\widehat U(\cdot;\varepsilon)\) admits analytic continuation to
		      \(\Omega_\varepsilon\). Let \(\Gamma_\varepsilon:=\partial\Omega_\varepsilon\),
		      oriented positively. Then \(\Gamma_\varepsilon\) is a simple closed contour
		      winding once around \(0\), the closed disk
		      \[
			      \overline{D\bigl(0,\rho(\varepsilon)^s\bigr)}
			      \subset \Omega_\varepsilon,
			      \qquad
			      \rho(\varepsilon):=\frac{1+\rho_*(\varepsilon)}{2},
		      \]
			  and \(z_*(\varepsilon)\notin \Gamma_\varepsilon\).
		      Moreover,
		      \[
			      M:=\sup_{0<\varepsilon<\varepsilon_0}\ \sup_{z\in\Gamma_\varepsilon}\ |\widehat U(z;\varepsilon)|<\infty.
		      \]
		\item\label{U2} \emph{Uniform Puiseux expansion to order $3/2$.}
		        For all $z\in\Delta(z_*(\varepsilon),\phi)$ with $|\delta_z|<\eta_0$, where
		      \[
			     \delta_z:=1-\frac{z}{z_*(\varepsilon)},
		      \]
		      one has a Puiseux expansion on the Taylor sheet of the form
		      \begin{equation}\label{eq:puiseux-uniform}
			      \widehat U(z;\varepsilon)
			      =
			      \lambda(\varepsilon)
			      +
				     s^{-1/2}\kappa(\varepsilon)\,\delta_z^{1/2}
			      +
				     c(\varepsilon)\,\delta_z
			      +
				    d(\varepsilon)\,\delta_z^{3/2}
			      +
			      E(z;\varepsilon),
		      \end{equation}
			    where $\delta_z^{1/2}$ is taken with $\arg\delta_z\in(-\pi,\pi)$, and the coefficients and remainder satisfy
		      \[
			      |\lambda(\varepsilon)|\ge c_0,\qquad |\kappa(\varepsilon)|\ge c_0,\qquad
			      |\lambda(\varepsilon)|+|\kappa(\varepsilon)|+|c(\varepsilon)|+|d(\varepsilon)|\le C_0,
		      \]
		      and $|E(z;\varepsilon)|\le C_0\,|\delta_z|^{2}$.
			  In addition, for each $p\ge1$, let $\mathcal R_p(\cdot;\varepsilon)$ denote the regular part in
		      the local expansion
		      \[
			      \widehat U(z;\varepsilon)^p
			      =
			      \lambda(\varepsilon)^p
			      +
			      b_p(\varepsilon)\,\delta_z^{1/2}
			      +
			      d_p(\varepsilon)\,\delta_z^{3/2}
			      +
			      \mathcal R_p(z;\varepsilon)
		      \]
		      on $\Omega_\varepsilon\cap\{|\delta_z|<\eta_0\}$, where $b_p(\varepsilon)$ and
		      $d_p(\varepsilon)$ are the corresponding Puiseux coefficients. We assume that each
		      $\mathcal R_p(\cdot;\varepsilon)$ extends holomorphically to the full disk
		      \[
			      |z|<R(\varepsilon)=|z_*(\varepsilon)|+\eta_0/2,
		      \]
		      and satisfies the uniform boundary bound
		      \[
			      \sup_{0<\varepsilon<\varepsilon_0}\ \sup_{|z|=R(\varepsilon)}
			      |\mathcal R_p(z;\varepsilon)|
			      \le B_0^{\,p}
			      \qquad (p\ge1).
		      \]
		      Here $\kappa(\varepsilon)$ is the square-root coefficient in the
		      original $x$-variable, so that
		      \[
			      U(x;\zeta(\varepsilon))
			      =
			      \lambda(\varepsilon)
			      +
			      \kappa(\varepsilon)\sqrt{1-x/x_*(\varepsilon)}
			      +
			      O\!\left(1-x/x_*(\varepsilon)\right).
		      \]
	\end{enumerate}
\end{definition}

\begin{lemma}[Transfer from a dominant simple $s$--orbit]\label{lem:coeff-asymp-dominant}
	Fix $\zeta$ and assume Definition~\ref{def:dominant-orbit}. Let $x_*(\zeta)$ be a representative
	dominant singularity so $|x_*(\zeta)|=\rho_*(\zeta)$, and assume that $U(\cdot;\zeta)$ is
	$\Delta$--analytic at each orbit point $\{e^{2\pi i j/s}x_*(\zeta)\}_{j=0}^{s-1}$ and that the
	representative point is \emph{simple} on the Taylor sheet in the sense of Proposition~\ref{prop:xi-star-system} (so the nondegeneracy conditions \eqref{eq:Fx-nondeg} and \eqref{eq:Fyy-nondeg} hold for the corresponding branch value $\lambda$ at $x_*$).
	Then for every $p\ge1$, writing
	\[
		U(x;\zeta)^p=\sum_{m\ge0} R_p(m;\zeta)\,x^{ms},
	\]
	one has the asymptotics
	\begin{equation}\label{eq:Rp-asymp}
		R_p(m;\zeta)=A_p(\zeta)\,x_*(\zeta)^{-ms}\,m^{-3/2}\Bigl(1+O(m^{-1})\Bigr),
		\qquad m\to\infty,
	\end{equation}
	and there exists $C>0$ (locally in $\zeta$) such that
	\[
		|A_p(\zeta)|\le C\,p\,|\lambda|^{p-1}.
	\]
\end{lemma}

	\begin{equation}
		R_p(m;\zeta)=A_p(\zeta)\,x_*(\zeta)^{-ms}\,m^{-3/2}\Bigl(1+O(m^{-1})\Bigr),
		\qquad m\to\infty,
	\end{equation}

\begin{lemma}[Critical limit of the transfer amplitudes]\label{lem:critical-coeff-asymp}
	Let $\zeta(\varepsilon)$ be a $C^1$ path with $\zeta(0)=\zeta_c\in\mathcal C$ and $\varepsilon\downarrow0$.
	Assume Definition~\ref{def:dominant-orbit} along $\zeta(\varepsilon)$ and that the dominant orbit consists
	of simple branch points on the Taylor sheet. Choose a representative $x_*(\varepsilon)$ and let
	$\lambda(\varepsilon)$ be the corresponding branch value defined in \eqref{eq:lambda-def}.
	Then there exist limits
	\[
		x_*(\varepsilon)\to x_{*,c},\qquad
		\lambda(\varepsilon)\to \lambda_c,\qquad
		\kappa(\varepsilon)\to\kappa_c\neq0,
	\]
	where $\kappa(\varepsilon)$ is the square--root coefficient in the Puiseux expansion of the Taylor sheet
	at $x=x_*(\varepsilon)$. In the variable \(z=x^s\) from
	Definition~\ref{def:uniform-delta}, the coefficient of
	\(\sqrt{1-z/z_*(\varepsilon)}\) is therefore
	\(s^{-1/2}\kappa(\varepsilon)\), up to the fixed choice of square-root branch.
	Moreover, for each fixed $p\ge1$ the amplitudes $A_p(\varepsilon)$ from Lemma~\ref{lem:coeff-asymp-dominant}
	satisfy
	\[
		A_p(\varepsilon)\to A_p^{(c)}\qquad (\varepsilon\downarrow0),
	\]
	with the explicit limit
	\begin{equation}\label{eq:Ap-critical-explicit}
	A_p^{(c)}
	=
	-\frac{s^{-1/2}}{2\sqrt{\pi}}\;p\,\kappa_c\,\lambda_c^{\,p-1},
	\end{equation}
	again up to the fixed choice of square-root branch. Finally, there exist $\varepsilon_0>0$ and
	$C>0$ such that for all $\varepsilon\in(0,\varepsilon_0)$ and all $p\ge1$,
	\[
		|A_p(\varepsilon)|\le C\,p\,|\lambda(\varepsilon)|^{p-1}.
	\]
\end{lemma}

\begin{proof}[Proof of Lemma~\ref{lem:critical-coeff-asymp}]
	Abbreviate $\zeta=\zeta(\varepsilon)$. By
	Lemma~\ref{lem:local-branch-data}, after shrinking the neighborhood if
	necessary there is a unique $C^1$ continuation
	$(\lambda(\varepsilon),x_*(\varepsilon))$ of the chosen representative
	branch-point parameters for $\varepsilon$ small.
	In particular,
	\[
		x_*(\varepsilon)\to x_{*,c}:=x_*(0),\qquad \lambda(\varepsilon)\to\lambda_c:=\lambda(0).
	\]
	Moreover, by Proposition~\ref{prop:xi-star-system} each $x_*(\varepsilon)$ is a simple branch point of the Taylor sheet,
	hence the Taylor branch has a Puiseux expansion at $x=x_*(\varepsilon)$ of the form
	\begin{equation}\label{eq:U-local-eps-proof}
		U(x;\zeta(\varepsilon))
		=
		\lambda(\varepsilon)+\kappa(\varepsilon)\sqrt{1-x/x_*(\varepsilon)}
		+O\!\left(1-x/x_*(\varepsilon)\right),
	\end{equation}
	For a simple branch point one may choose the square-root branch so that
	\[
		\kappa(\varepsilon)^2=
		\frac{2x_*(\varepsilon)\,\partial_xF(\lambda(\varepsilon),x_*(\varepsilon);\zeta(\varepsilon))}
		{\partial_y^2F(\lambda(\varepsilon),x_*(\varepsilon);\zeta(\varepsilon))}.
	\]
	The right-hand side depends continuously on $\varepsilon$ and stays nonzero, so after fixing one branch sign
	consistently we obtain $\kappa(\varepsilon)\to\kappa_c\neq0$.
	By the change of variable $z=x^s$ one has
	\[
		1-\frac{z}{z_*(\varepsilon)}
		=
		1-\Bigl(\frac{x}{x_*(\varepsilon)}\Bigr)^s
		=
		s\Bigl(1-\frac{x}{x_*(\varepsilon)}\Bigr)
		+
		O\!\left(\Bigl(1-\frac{x}{x_*(\varepsilon)}\Bigr)^2\right),
	\]
	so the coefficient of $\sqrt{1-z/z_*(\varepsilon)}$ is
	$s^{-1/2}\kappa(\varepsilon)$, up to the fixed choice of square-root branch.

	Fix $p\ge1$. Lemma~\ref{lem:coeff-asymp-dominant} applied at $\zeta(\varepsilon)$ yields
	\eqref{eq:Rp-asymp} with amplitude $A_p(\varepsilon)$, and the uniform estimate
	$|A_p(\varepsilon)|\le C\,p\,|\lambda(\varepsilon)|^{p-1}$.

	Using the explicit amplitude formula from the proof of Lemma~\ref{lem:coeff-asymp-dominant}, we obtain
	\[
		A_p(\varepsilon)
		=
		-\frac{s^{-1/2}}{2\sqrt{\pi}}\;
		p\,\kappa(\varepsilon)\,\lambda(\varepsilon)^{p-1}.
	\]
	Passing to the limit gives $A_p(\varepsilon)\to A_p^{(c)}$ with
	\[
		A_p^{(c)}
		=
		-\frac{s^{-1/2}}{2\sqrt{\pi}}\;
		p\,\kappa_c\,\lambda_c^{\,p-1},
	\]
	up to the conventional choice of the square--root branch.
\end{proof}

\end{document}
